\section{Semantic interfaces and the certified criterion}
\label{sec:semantics}

The reference question is whether a merged implementation state has an
explanation in an independently defined sequential language, using exactly
its original events. Agreement of implementation replays alone does not answer
this question. The definitions below describe the checked concrete and
invariant-scoped interfaces. They distinguish the globally guarded sufficient
class from the certified, represented-prefix class. An implementation need
not inhabit the former to have a certificate in the latter.

\begin{definition}[Events, effects, and payload policy]
\label{def:events}
An implementation has a concrete state space $\Sigma$, initial state
$\sigma_0$, operations $O$, queries $Q$, answers $V$, a total effector
$u:\Sigma\times(\mathbb N\times\mathbb N\times O)\to\Sigma$, a query
$\operatorname{read}:\Sigma\times Q\to V$, and a three-way merge
$m:\Sigma^3\to\Sigma$. An event is $e=(t,r,o)$, with projections
$e.t,e.r,e.o$. For a list $\pi$, $u^*(s,\pi)$ is its left fold from $s$.
A payload policy is a relation $P\subseteq O\times O$; write
$a\mathrel{P}b$ for $P(a.o,b.o)$. It does not inspect the events' external
timestamps or replica identifiers.
A replay context supplies eligible events $\mathcal E_C$ and visibility
$\mathord{\prec_C}$. Set
\[
 a\parallel_C b\ \Longleftrightarrow\
 \neg(a\prec_C b)\land\neg(b\prec_C a),\qquad
 \operatorname{distinct}(a,b)\ \Longleftrightarrow\ a.t\ne b.t.
\]
Concrete raw commutation is
\[
 a\mathrel{\bowtie}b\quad\Longleftrightarrow\quad
 \forall s\in\Sigma,\quad u(u(s,a),b)=u(u(s,b),a).
\]
\end{definition}
These are relations on full events, even though $P$ acts on payloads.

\begin{definition}[Globally guarded replay laws]
\label{def:guards}
The global guarded sufficient class requires the following, for full events:
\begin{align*}
&a.t\ne b.t\ \land\ a.r\ne b.r
 \quad\Longrightarrow\quad
 (\neg(a\bowtie b)\ \Longleftrightarrow\ aPb\lor bPa),\\
&a.t\ne b.t\ \land\ b.t\ne c.t
 \quad\Longrightarrow\quad \neg(aPb\land bPc).
\end{align*}
For every $s\in\Sigma$, intervening event list $\beta$, and $a,b,c$ with
pairwise distinct timestamps, it also requires
\[
 aPb\land\neg(b\bowtie c)\quad\Longrightarrow\quad
 u\bigl(u^*(u(u(s,b),a),\beta),c\bigr)
 =u\bigl(u^*(u(u(s,a),b),\beta),c\bigr).
\]
There is no different-replica guard on the last equation. Its absorber premise
is actual noncommutation, rather than the policy edge $bPc$.
\end{definition}
The no-chain clause allows $a=c$ when the two adjacent timestamp guards hold;
therefore a payload self-edge is not rescued by external event freshness.
These global equations quantify over arbitrary concrete scratch states and
intervening lists. They are not the scoped laws below.

\begin{definition}[Invariant-scoped concrete commutation and closure]
\label{def:invariant}
For a fixed concrete state predicate $I:\Sigma\to\mathrm{Prop}$, define
\[
 a\mathrel{\bowtie_I}b\quad\Longleftrightarrow\quad
 \forall s\in\Sigma,\ I(s)\Longrightarrow
 u(u(s,a),b)=u(u(s,b),a).
\]
An invariant closure certificate at $C$ supplies
\begin{align*}
&I(\sigma_0),\\
&\forall s,e,\quad I(s)\land e\in\mathcal E_C\Longrightarrow I(u(s,e)),\\
&\forall l,a,b,\quad I(l)\land I(a)\land I(b)
 \Longrightarrow I(m(l,a,b)).
\end{align*}
\end{definition}
Commutation here has no simultaneous-readiness or reissuability premise.
Closure includes replaying an already issued event out of causal order and
repeating it. A port must independently justify the predicate and its closure;
it cannot exclude a reachable state merely because that state refutes an equation.

\begin{definition}[Version-local invariant order]
\label{def:order}
For $E\subseteq\mathcal E_C$, use the \emph{same} predicate $I$ in both
noncommutation tests:
\begin{align*}
 a\mathrel{\operatorname{lo}^{I}_{C,E}}b\quad\Longleftrightarrow\quad
 &\bigl(a\prec_C b\land\neg(a\bowtie_I b)\bigr)\\
 {}\lor{}&\bigl(a\parallel_C b\land aPb\land
 \neg\exists c\in E.\ b\prec_C c\land\neg(b\bowtie_I c)\bigr).
\end{align*}
Only events in $E$ can suppress a concurrent policy edge as absorbers.
For $I(s)\equiv\mathrm{True}$ this is the raw order
$\operatorname{lo}_{C,E}$, using $\bowtie$ throughout.
A list enumerates $E$ when it has no duplicates and its membership is exactly
$E$. It respects a relation when every related pair in the list appears in
that relation's direction.
\end{definition}
Neither order is the full visibility relation. In particular,
causal implementation-commuting events may be reordered by this order.
Readiness adds causal constraints to implementation replay; it does not
silently redefine the public semantic order.

\begin{definition}[Original issuance, certified executions, and readiness]
\label{def:issuance}
An issuance discipline is a predicate $G(e,s)$ checked at the materialized
replica head from which $e$ was minted. Mint honesty at $C$ requires, for every
$e\in\mathcal E_C$, a list $\pi_e$ satisfying
\[
 \pi_e\text{ enumerates }\{p\in\mathcal E_C\mid p\prec_C e\},\qquad
 \pi_e\text{ respects }\prec_C,\qquad G(e,u^*(\sigma_0,\pi_e)).
\]
An issued application is an actual raw store application with its original
head/version evidence and $G(e,s)$; non-application steps preserve the raw
fork/merge rule. A certified execution starts at the initial store and, at
every transition, retains mint honesty before and after the issued step.
The widened execution permits recursively synthesized virtual merge bases.
It has its own certified reachability relation, with the ordinary steps as its
base fragment.

A replay scope is $(C,F,R)$, where $F$ is its fixed eligible event set and
$R(H,s)$ records representation of the already replayed prefix $H$. Define
\[
 \operatorname{Ready}(H,e)\quad\Longleftrightarrow\quad
 e\in F\land e\notin H\land
 \forall p\in F,\ p\prec_C e\Longrightarrow p\in H.
\]
A replay is legal from $H$ by the inductive rules
\[
 \operatorname{Legal}(H,[]),\qquad
 \frac{\operatorname{Ready}(H,e)\quad
       \operatorname{Legal}(H\cup\{e\},\pi)}
      {\operatorname{Legal}(H,e::\pi)}.
\]
\end{definition}
The issuance guard is retained as evidence about the original mint. It is not
checked again when the effector is replayed at a different scratch state.

\begin{definition}[Scoped replay and policy obligations]
\label{def:scoped-laws}
The invariant replay laws require
\[
 R(H,s)\Longrightarrow I(s),\qquad
 R(H,s)\land\operatorname{Ready}(H,e)
 \Longrightarrow R(H\cup\{e\},u(s,e)).
\]
They also require the concrete diamond
$u(u(s,a),b)=u(u(s,b),a)$ whenever $R(H,s)$, both events are ready,
$a\parallel_C b$, and neither direction of
$\operatorname{lo}^{I}_{C,F}$ orders them.
The accompanying policy laws require
\[
 \neg(a\bowtie_I b)\Longleftrightarrow aPb\lor bPa
\]
only for $a,b\in F$ that are concurrent, have distinct timestamps, and come
from different replicas. No-chain is required for $a,b,c\in F$ with
$a.t\ne b.t$ and $b.t\ne c.t$, as in Definition~\ref{def:guards}.
The package consists of both replay laws and policy laws.
\end{definition}
An alternative represented-prefix commutation test instead quantifies over represented
prefixes where both events are ready. That definition is distinct from
$\bowtie_I$; replacing the latter by readiness would make causal-pair tests
potentially vacuous.

\begin{theorem}[Scoped implementation replay equality]
\label{thm:scoped-replay}
Suppose the replay laws of Definition~\ref{def:scoped-laws} hold,
$R(H,s)$, and $\pi_1,\pi_2$ enumerate the same event set $A$. If both lists are
legal from $H$ and each respects both $\prec_C$ and
$\operatorname{lo}^{I}_{C,F}$, then
\[
 u^*(s,\pi_1)=u^*(s,\pi_2).
\]
\end{theorem}
Its hypotheses are scoped replay laws, not the global class of
Definition~\ref{def:guards}. This theorem proves concrete implementation
agreement. Independent sequential legality is a further obligation.

\begin{definition}[Independent history language and contextual conflict]
\label{def:history}
For an explicitly chosen update alphabet $U$, let labels be
$\operatorname{upd}(x)$, $x\in U$, and $\operatorname{qry}(q,v)$.
A history specification $S$ is a prefix-closed set of finite label lists
containing the empty list. A labelled transition system defines such a language
by accepting every finite run from its initial state; no determinism or totality
is assumed. Thus acceptance is a proposition, not a uniquely determined
resulting state.
Define contextual specification commutation by
\begin{align*}
 x\mathrel{\bowtie_S}y\quad\Longleftrightarrow\quad
 \forall\rho,\tau,\quad
 &\rho\cdot[\operatorname{upd}(x),\operatorname{upd}(y)]\cdot\tau\in S\\
 &\Longleftrightarrow
 \rho\cdot[\operatorname{upd}(y),\operatorname{upd}(x)]\cdot\tau\in S.
\end{align*}
For a fixed label projection $f:(\mathbb N\times\mathbb N\times O)\to U$,
let $\operatorname{labels}_f(\pi)$ retain the original event order while
mapping each update through $f$. The independent specification visibility
relation is
\[
 a\mathrel{\operatorname{sv}^{f,S}_C}b
 \quad\Longleftrightarrow\quad
 a\prec_C b\land\neg\bigl(f(a)\bowtie_S f(b)\bigr).
\]
\end{definition}
Full-input specifications use $f=\mathrm{id}$.
Payload projection $f(e)=e.o$ is an explicit interface choice. It cannot be
assumed adequate for specifications that distinguish original enqueue or
insertion identities.

\begin{definition}[The full-event invariant criterion]
\label{def:ra}
For every allocated version $v$ with stored record $(s,E)$ and every query $q$,
there must exist \emph{one} list $\pi$ such that
\begin{enumerate}[nosep]
\item $\pi$ enumerates exactly $E$;
\item $\pi$ respects $\operatorname{lo}^{I}_{C,E}$;
\item $\pi$ respects $\operatorname{sv}^{\mathrm{id},S}_C$;
\item $\operatorname{labels}_{\mathrm{id}}(\pi)\cdot
 [\operatorname{qry}(q,\operatorname{read}(s,q))]\in S$.
\end{enumerate}
\end{definition}
The universal invariant $I(s)\equiv\mathrm{True}$ recovers the raw
full-event criterion, including specification visibility. Checking every
allocated version implies the corresponding guarantee for current replica
heads. A projected criterion replaces the specification alphabet through $f$;
it still preserves the specification conflicts of those projected labels.
The manuscript's literal criterion omits item~3 and uses payload labels.
These are distinct criteria, rather than alternate notations for the same
proposition.

An independent sequential simulation may use a relation, or an abstraction
function $\alpha:\Sigma\to\Phi$, to relate implementation states to the
\emph{specification's} states. This does not replace concrete equality in the
merge VCs, diamonds, or Definition~\ref{def:invariant}. Universal fold soundness
is one sufficient bridge. Certified Queue and RGA results instead establish
appropriate original-event witnesses for represented histories. Neither
implementation replay equality nor existence of an implementation fold alone
establishes item~4, and a specification-legal witness must also satisfy item~3.
